\section{The Hook-Components conjecture and reconstruction}
\label{sec:augmented}

The sign representation turns an exterior degree into a hook label.
To use this observation on coinvariants, we show that every
positive odd trace already vanishes on the sign component once the
even traces have been removed.  The argument uses multiplication
and augmentation alone, so we work at the level of an arbitrary
pointed algebra.

This section works over any commutative unital ring $K$ in which $n!$ is
invertible.  Let $A$ be a commutative unital $K$-algebra
with an augmentation $\eps:A\to K$, and let
$\mathfrak m=\ker\eps$.  Set
\[
 B=A^{\otimes n},\qquad
 \eps_B=\eps^{\otimes n}.
\]
The symmetric group permutes the tensor factors.
For $a\in A$, let $a^{(i)}$ be $a$ in the $i$th tensor
factor and $1$ in all other factors.
Let $N_n$ be the permutation representation with basis
$\theta_1,\ldots,\theta_n$, and set
\[
 S=B\otimes\Lambda N_n.
\]
Its augmentation $\eps_S$ extends $\eps_B$ and sends
every $\theta_i$ to zero.

\begin{definition}\label{def:augmented-coinvariants}
The ordinary and exterior coinvariant quotients at the
augmentation are
\[
 I_A=B\,\ker(\eps_B|_{B^{S_n}}),\qquad R_A=B/I_A,
\]
and
\[
 J_A=S\,\ker(\eps_S|_{S^{S_n}}),\qquad Q_A=S/J_A.
\]
\end{definition}

If $A$ is connected and graded, these are the usual
quotients by invariants with zero constant term.
The augmentation definition also makes sense for
ungraded and nonreduced algebras.
Put $\theta_\Sigma=\sum_i\theta_i$.
As $n$ is invertible, $N_n/K\theta_\Sigma$ identifies
canonically with the standard representation $W_n$.

\begin{theorem}[Augmented sign transfer]\label{thm:augmented-transfer}
The natural quotient induces an isomorphism
\begin{equation}\label{eq:augmented-transfer}
 (R_A\otimes\Lambda W_n)^{\sgn}
       \xrightarrow{\ \sim\ }Q_A^{\sgn}.
\end{equation}
It preserves exterior degree, is natural under augmented
algebra homomorphisms, commutes with arbitrary base change, and
commutes with all augmented
automorphisms of $A$.  It also preserves any additional
grading respected by the augmentation.
\end{theorem}

For a polynomial algebra in two generators, this is the
sign-transfer theorem of Gonz\'alez, Lentfer and Mularczyk
\cite[Theorem~4.21]{GLM}.  Their complementary-averaging and
determinant argument \cite[Lemmas~4.13--4.15]{GLM} extends to the
present setting because collisions multiply elements of the
algebra and positive traces remain in its augmentation ideal.

\subsection{Trace generators}

For $a\in A$ and $T\subseteq[n]$, write
\[
 t_a(T)=\sum_{i\in T}a^{(i)},\qquad
 q_a=\sum_{i=1}^n\theta_i a^{(i)},\qquad
 t_a=t_a([n]).
\]

\begin{lemma}\label{lem:trace-generation}
The algebra $B^{S_n}$ is generated by $1$ and $t_a$
for $a\in\mathfrak m$.
The superalgebra $S^{S_n}$ is generated by these even
elements together with $q_a$ for $a\in A$.
Consequently
\begin{equation}\label{eq:trace-ideals}
 I_A=\langle t_a:a\in\mathfrak m\rangle,\qquad
 J_A=\langle t_a:a\in\mathfrak m;\
                    q_a:a\in A\rangle.
\end{equation}
\end{lemma}

\begin{proof}
Use the split $K$-module decomposition $A=K1\oplus\mathfrak m$.
Pure tensors in $B$, together with exterior monomials,
span $S$ as a $K$-module; averaging spans its invariant submodule.
An averaged tensor can be described by finitely many
occupied labels, each carrying an even decoration from
$\mathfrak m$ or an odd decoration $\theta a$ with
$a\in A$.

Multiply the traces corresponding to those decorations.
The terms with pairwise distinct labels give the desired
orbit sum, up to the multiplicity factorials of repeated
even decorations.  These factorials divide $n!$.
A repeated odd decoration has zero distinct-label
antisymmetrization, so no division by a vanishing scalar
is required.  Every other term has a collision and uses
fewer occupied labels.  Two odd decorations at one label
give zero.  An even collision is multiplication in $A$,
and a collision involving an even decoration from
$\mathfrak m$ stays in $\mathfrak m$.
Induction on the number of occupied labels proves the
two generation statements.  All generators except $1$
have zero augmentation, giving \eqref{eq:trace-ideals}.
\end{proof}

For a set $T$ of labels, let $\Av_T$ and $\Alt_T$ denote
the normalized symmetric and alternating averaging
operators over $S_T$.
The averaging identity below proves sign transfer. A trace relation
moves a factor from the complementary labels onto an
alternating block, where the resulting determinant has a repeated
row and vanishes.

\begin{lemma}[Complementary averaging]\label{lem:complement-average}
Let $D\subseteq[n]$, let $i\notin D$, and put
$C=D\cup\{i\}$.  Suppose $f\in B$ is alternating on
$C$ and symmetric on $C^c$.
For every $a\in\mathfrak m$,
\begin{equation}\label{eq:complement-average}
 \Av_{D^c}\bigl(a^{(i)}f\bigr)\in I_A.
\end{equation}
\end{lemma}

\begin{proof}
Separate the tensor factors on $C$ and $C^c$.
The symmetric and alternating averaging operators are idempotents
on separate tensor factors.  Their images are split direct summands;
tensoring these splittings identifies the module with the indicated
symmetries as the tensor product of the alternating module on $C$
and the invariant module on $C^c$.  No flatness is required.
By Lemma~\ref{lem:trace-generation}, the latter is
generated by the traces $t_b(C^c)$ with $b\in\mathfrak m$.
Modulo $I_A$,
\[
 t_b(C^c)=-t_b(C).
\]
After making these substitutions, $f$ is congruent to
a sum of elements using only the tensor factors on
$C$ and alternating on $C$.  Multiplication by $a^{(i)}$
and averaging preserve congruence modulo $I_A$.

Write $c_1,\ldots,c_h$ for the labels of $C$, with
$c_h=i$.
Alternating tensors on these factors are spanned as a $K$-module by
determinants
\[
 F=\det\bigl(b_j^{(c_\ell)}\bigr)_{1\leq\ell,j\leq h},
        \qquad b_j\in A.
\]
Indeed, the alternating projector sends a pure tensor to its
determinant divided by $h!$.  Thus it suffices to prove \eqref{eq:complement-average}
for $F$.
Let $M_j$ be the cofactor of the entry in its last row
and $j$th column.  For $r\in D^c$, let $F_r$ be the
same determinant with its last row evaluated at label
$r$.  The cofactor $M_j$ uses labels in $D$ only.
It follows that
\begin{align*}
 |D^c|\,\Av_{D^c}(a^{(i)}F)
   &=\sum_{r\in D^c}a^{(r)}F_r\\
   &=\sum_{j=1}^h M_j\,t_{ab_j}(D^c)\\
   &\equiv-\sum_{r\in D}a^{(r)}
                     \sum_{j=1}^h M_j b_j^{(r)}
                          \pmod{I_A}\\
   &=0.
\end{align*}
The congruence uses $ab_j\in\mathfrak m$.
Each final inner sum is a determinant with a repeated
row.  Since $|D^c|$ is invertible, the assertion follows.
\end{proof}

\subsection{Odd invariants on the sign component}

\begin{lemma}\label{lem:odd-annihilation}
For $a\in\mathfrak m$, multiplication by $q_a$
annihilates the sign submodule of
$R_A\otimes\Lambda N_n$.
\end{lemma}

\begin{proof}
Sign tensors are spanned by the images of
\[
 \Alt_{[n]}(\theta_J f),
 \qquad
 f\text{ symmetric on }J
       \text{ and alternating on }J^c,
\]
where $\theta_J$ is an exterior monomial with support $J$.
Indeed, average first over the stabilizer
$S_J\times S_{J^c}$: the sign of the exterior monomial
cancels the sign projector on $S_J$.

Since $q_a$ is invariant, its product with such a tensor
is a sum over $i\in J^c$ of terms
\[
 \Alt_{[n]}(\theta_i\theta_J\,a^{(i)}f).
\]
Fix one term and put $D=J^c\setminus\{i\}$.
The exterior factor has support $D^c=J\cup\{i\}$.
Inside the global sign projector, its coefficient can
be replaced by its symmetric average over $D^c$:
the exterior sign again cancels the sign projector.
But $f$ is alternating on $D\cup\{i\}$ and symmetric
on its complement.  Lemma~\ref{lem:complement-average}
therefore says that the averaged coefficient
$\Av_{D^c}(a^{(i)}f)$ is zero in $R_A$.
Every term vanishes.
\end{proof}

\begin{proof}[Proof of Theorem~\ref{thm:augmented-transfer}]
First quotient $S$ by the even generators $t_a$,
$a\in\mathfrak m$.  The result is
$R_A\otimes\Lambda N_n$.
Next quotient by $q_1=\theta_\Sigma$; the result is
$R_A\otimes\Lambda W_n$.
Since $A=K1\oplus\mathfrak m$, the only generators
of $J_A$ still to impose are $q_a$ with
$a\in\mathfrak m$.

Lemma~\ref{lem:odd-annihilation} kills these operators
on the sign part before imposing $q_1$.  Exactness of
the sign projector shows that they also kill the sign
part after imposing $q_1$.
Moreover multiplication by each $q_a$ is
$S_n$-equivariant.  The sign part of its image is
therefore its image on the sign part, and is zero.
The remaining ideal has no sign component.
Quotienting by that ideal preserves the sign submodule
and proves \eqref{eq:augmented-transfer}.
All the constructions used the canonical quotient
maps and averaging projectors, which also proves
the asserted naturality and grading compatibility.

For a homomorphism $K\to K'$, put $A'=A\otimes_KK'$.
The averaging and sign projectors are split idempotents, so their
images commute with arbitrary scalar extension.  The augmentation
on each invariant algebra is split by the unit; hence its kernel
also commutes with scalar extension.  Present $I_A$ and $J_A$ as
images of the multiplication maps from these kernels.  Right
exactness of tensor product then gives canonical isomorphisms
\[
 R_A\otimes_KK'\cong R_{A'},\qquad
 Q_A\otimes_KK'\cong Q_{A'}.
\]
The splitting $N_n=K\theta_\Sigma\oplus W_n$ and its finite-free
exterior powers commute with scalar extension as well.  Thus so
does \eqref{eq:augmented-transfer}.
\end{proof}

\begin{corollary}[Hook extraction on a pointed algebra]
\label{cor:pointed-hooks}
For $0\leq d\leq n-1$ there are natural isomorphisms
\[
 \Hom_{K[S_n]}(\sgn,Q_A[d])
       \cong
 \Hom_{K[S_n]}(S^{(d+1,1^{n-d-1})},R_A),
\]
where $[d]$ denotes exterior degree and $S^\mu$ is the base
change to $K$ of the integral Specht module.
They commute with arbitrary base change and every augmented
automorphism of $A$.
\end{corollary}

\begin{proof}
Theorem~\ref{thm:augmented-transfer} identifies the
left side with
$\Hom_{K[S_n]}(\sgn,R_A\otimes\Lambda^dW_n)$.
The standard pairing on $N_n$ restricts to a perfect pairing on
$W_n$: the complementary vector $\theta_\Sigma$ has norm $n$,
a unit.  Since $W_n$ is finite free and $\det W_n=\sgn$,
exterior duality and the integral hook Specht model
\cite[Section~5.5(b)]{DanzHofmann} give
\[
 \sgn\otimes\Lambda^dW_n
       \cong S^{(d+1,1^{n-d-1})}.
\]
Finite-free tensor-Hom adjunction gives the result.  Equivariant
Hom is the image of an averaging projector on ordinary Hom, so
these identifications also commute with arbitrary base change.
\end{proof}

\begin{example}
Let $A=K[z_1,\ldots,z_k]/L$, where every element of
$L$ vanishes at the origin.  Evaluation at the origin
defines an augmentation.
Theorem~\ref{thm:augmented-transfer} applies whether or
not the pointed affine scheme $\operatorname{Spec}A$
is smooth or reduced.  It applies, for example, to
$K[u,v]/(uv)$ and to $K[u_1,\ldots,u_k]/(u_1,\ldots,u_k)^m$.
The theorem also applies over $\Z[1/n!]$ and to every square-zero
extension $A=K\oplus M$, with $M^2=0$, for an arbitrary $K$-module
$M$.  In particular, $M$ need not be flat.  Over a field of
characteristic $p>n$, the same theorem holds.
The characteristic-zero Schur character identities in
the following subsections use the polynomial case and
additional representation-theoretic input.
\end{example}


\subsection{Supersymmetric specialization}
\label{sec:hook-identity}

Return to a field $K$ of characteristic zero and take
$A=K[z_1,\ldots,z_k]$ with augmentation at the origin.
Then $R_A=R_{n,k}$.
The algebra $Q_A$ has one exterior permutation alphabet in
addition to these bosonic alphabets.
For a row alphabet $\bq=(q_1,\ldots,q_k)$ and one odd
variable $u$, use the positive supersymmetric convention
\begin{equation}\label{eq:superschur}
 s_\lambda(\bq\mid u)
   =\sum_{\lambda/\rho\text{ a vertical strip}}
           s_\rho(\bq)\,u^{|\lambda|-|\rho|}.
\end{equation}
In particular,
$[u^d]s_\lambda(\bq\mid u)
     =(e_d^\perp s_\lambda)(\bq)$.

Lentfer's diagonal supersymmetry theorem
\cite[Theorem~1.2 and Corollary~1.3]{Lentfer}
states that the same universal coefficients
$c_{\lambda,\mu}^{(n)}$ describe the full coinvariant
characters in arbitrary bosonic and fermionic alphabets.
For the sign component with one odd alphabet, it says
\begin{equation}\label{eq:lentfer-sign}
 \ch\Hom_{S_n}(\sgn,Q_A)
       =\sum_\lambda
          c_{\lambda,(1^n)}^{(n)}s_\lambda(\bq\mid u).
\end{equation}

\begin{theorem}[All full hook identities]\label{thm:full-hook}
For $0\leq d\leq n-1$,
\begin{equation}\label{eq:full-hook}
 C_{(d+1,1^{n-d-1})}^{(n)}=e_d^\perp A_n.
\end{equation}
Equivalently,
\[
 c_{\rho,(d+1,1^{n-d-1})}^{(n)}
   =\sum_{\lambda/\rho\text{ a vertical }d\text{-strip}}
          c_{\lambda,(1^n)}^{(n)}.
\]
\end{theorem}

\begin{proof}
By Corollary~\ref{cor:pointed-hooks}, the coefficient
of $u^d$ on the left side of \eqref{eq:lentfer-sign}
is $C_{(d+1,1^{n-d-1})}^{(n)}(\bq)$.
By \eqref{eq:superschur}, the coefficient on the
right side is $(e_d^\perp A_n)(\bq)$.
The equality holds for every $k$, and hence holds
as an identity of stable symmetric functions.
\end{proof}

The full hook identity follows from augmented sign transfer and
diagonal supersymmetry.  The two-bosonic-alphabet transfer in
\cite{GLM} gives its two-variable specialization; the augmented
construction applies in every row dimension.  We next construct the
corresponding identity inside the operator module.

\begin{remark}[Order of skewing and specialization]
\label{rem:skew-specialization}
In \(k\) alphabets the formula is
\[
 C_{(d+1,1^{n-d-1})}^{(n)}(\bq)=(e_d^\perp A_n)(\bq).
\]
It is generally impossible to replace $A_n$ by the
sum of its terms of length at most $k$ before skewing.
A vertical strip can remove the last box in each of
$d$ rows, so source partitions of length $k+d$ can
contribute to a target of length $k$.
\end{remark}

\subsection{Hook transfer in the polarization module}
\label{sec:r2-operator-hooks}

Invariant differential forms give the operator counterpart.  An
odd differential trace can be realized by a polarization into a
fresh row.  Repeating this operation and alternating the fresh
rows lifts every hook vector to an alternating vector, which is
the required inverse to exteriorization.

We work over a field of characteristic zero.  For a finite row space
$V$, write $P(V)=\Sym(V\otimes N_n)$, where $N_n=K^n$ is the permutation
representation.  We use the polynomial row convention for the
variables.  All the subspaces below are polynomial row submodules
compatible with enlarging the row space, with taking homogeneous row
components, and with restriction obtained by setting auxiliary rows
to zero.  Let $\mathcal M(V)$
be a subspace of the diagonal harmonics, stable under constant
coefficient partial derivatives and generalized polarizations.
Its stable alternating character is denoted by $A^{\mathcal M}$.
For $\Harm$ these compatibilities follow directly from the
equations.  For $\Pol$ they follow from the commuting normal form
and its specialization statement,
Corollary~\ref{cor:r3-exact-specialization}.  

For a nonzero multiindex $\alpha$ on the old rows and a fresh row $z$,
the mixed polarization
\begin{equation}\label{eq:r2-mixed-polarization}
 E_{z,\alpha}=\sum_{i=1}^n z_i\partial_i^\alpha
\end{equation}
is an allowed operation.  Indeed, it is a linear combination of
conjugates of $\sum_i z_i\partial_{x_i}^{|\alpha|}$ under changes of
the old rows that fix $z$: pure powers of linear forms span the
homogeneous differential symbols of degree $|\alpha|$.

For $0\leq j\leq n-1$, add $j$ auxiliary rows $z_1,\ldots,z_j$.
Let $\mathscr E_j^{\mathcal M}(V)$ be the space of elements of
$\mathcal M(V\oplus K^j)^{\sgn}$ that are multilinear in these rows
and alternating under their permutations.  Such a polynomial has a
unique expression
\[
 F(X,z)=\sum_{i_1,\ldots,i_j}
       f_{i_1,\ldots,i_j}(X)z_{1,i_1}\cdots z_{j,i_j},
\]
with alternating coefficient tensor.  Set
\begin{equation}\label{eq:r2-exteriorization-map}
 \Psi_j(F)=\sum_{i_1,\ldots,i_j}
           f_{i_1,\ldots,i_j}(X)
                 \theta_{i_1}\cdots\theta_{i_j}.
\end{equation}
Here the $\theta_i$ generate $\Lambda N_n$.  Translation invariance in
each auxiliary row puts the coefficient tensor in $W_n^{\otimes j}$,
where $W_n=\ker(N_n\to K)$; hence the image is in $\Lambda^jW_n$.
Taking coefficients in the auxiliary rows is a sequence of partial
derivatives followed by setting those rows to zero, so every
$f_{i_1,\ldots,i_j}$ belongs to $\mathcal M(V)$.
Consequently
\begin{equation}\label{eq:r2-exteriorization-injection}
 \Psi_j:\mathscr E_j^{\mathcal M}(V)
       \lhook\joinrel\longrightarrow
          (\mathcal M(V)\otimes\Lambda^jW_n)^{\sgn}
\end{equation}
is a $\GL(V)$-equivariant injection.  Injectivity follows directly:
exteriorization is multiplication by $j!$ on alternating tensors.
The case $j=0$ is the identity.

Branching the stable alternating character from $V\oplus K^j$ to
$\GL(V)\times\GL_j$, and selecting the determinant representation of
$\GL_j$, gives
\begin{equation}\label{eq:r2-exteriorization-character}
 \ch\mathscr E_j^{\mathcal M}(V)
          =(e_j^\perp A^{\mathcal M})(V).
\end{equation}
Indeed, a multilinear alternating function of the auxiliary rows
transforms by their determinant, and the relevant
Littlewood--Richardson coefficients are
$c_{\rho,(1^j)}^\lambda$.

\begin{theorem}[Operator hook transfer]\label{thm:r2-operator-hook-transfer}
Let $\mathcal M$ be a compatible system of harmonic submodules generated
from alternating polynomials by constant coefficient derivatives and
generalized polarizations.  Then \eqref{eq:r2-exteriorization-injection}
is an isomorphism for every finite row space $V$ and every
$0\leq j\leq n-1$.  In particular,
\begin{equation}\label{eq:r2-operator-hook-character}
 C_{(j+1,1^{n-j-1})}^{\mathcal M}
                  =e_j^\perp A^{\mathcal M}.
\end{equation}
For $\mathcal M=\Pol$, this is Bergeron's Hook-Components Conjecture.
\end{theorem}

\begin{proof}
Put $\mathcal D=K[\partial_{x_i^{(a)}}]$.  By
Lemma~\ref{lem:r2-normal-order}, the evaluation map
\[
 \mathcal D\otimes\mathcal B(V)\longrightarrow\mathcal M(V),
 \qquad D\otimes g\longmapsto Dg,
\]
is surjective, and $S_n$ acts by its sign on $\mathcal B(V)$.
Exactness of averaging therefore shows that every element of the
right side of \eqref{eq:r2-exteriorization-injection} is a sum of
invariant differential forms in
$(\mathcal D\otimes\Lambda^jW_n)^{S_n}$ acting on elements of
$\mathcal B(V)$.

The invariant algebra $(\mathcal D\otimes\Lambda N_n)^{S_n}$ is generated by
the even traces and odd traces
\[
 p_\alpha=\sum_i\partial_i^\alpha
       \quad(|\alpha|>0),\qquad
 q_\alpha=\sum_i\theta_i\partial_i^\alpha
       \quad(|\alpha|\geq0).
\]
This is Lemma~\ref{lem:trace-generation}, applied to the polynomial
algebra of differential symbols.  Passing to $\Lambda W_n$ sets
$q_0=\sum_i\theta_i$ to zero.  Averaging is exact, so the images
of these generators still generate all invariants in the quotient.
The $p_\alpha$ commute with all the $q_\beta$ and annihilate the
harmonic polynomial on which the form acts.  A form of exterior
degree $j$ consequently reduces to a linear combination of
\begin{equation}\label{eq:r2-hook-odd-word}
             q_{\alpha_1}\cdots q_{\alpha_j}g,
       \qquad |\alpha_s|>0,\quad g\in\mathcal B(V).
\end{equation}

For this word, introduce the $j$ fresh rows and form
\[
 \widetilde F
       =E_{z_1,\alpha_1}\cdots E_{z_j,\alpha_j}g.
\]
All derivatives in these operators use only the old rows.
Thus the operators commute, their product is multilinear in the
fresh rows, and
\[
 \widetilde F
   =\sum_{i_1,\ldots,i_j}
        z_{1,i_1}\cdots z_{j,i_j}
        \partial_{i_1}^{\alpha_1}\cdots
        \partial_{i_j}^{\alpha_j}g.
\]
It belongs to $\mathcal M(V\oplus K^j)^{\sgn}$ by
\eqref{eq:r2-mixed-polarization}.  Let $F$ be its normalized
alternation under permutations of the fresh rows.
Then $F\in\mathscr E_j^{\mathcal M}(V)$.
Exteriorization already alternates these rows, so
\[
                  \Psi_j(F)
                       =q_{\alpha_1}\cdots q_{\alpha_j}g.
\]
Every expression \eqref{eq:r2-hook-odd-word} is therefore in the image
of $\Psi_j$, proving surjectivity.  The character identity follows from
\eqref{eq:r2-exteriorization-character} and
$\sgn\otimes\Lambda^jW_n\cong S^{(j+1,1^{n-j-1})}$.
\end{proof}

The construction also fixes the order of operations in
\eqref{eq:r2-operator-hook-character}: add the $j$ auxiliary rows,
select their exterior component, and then specialize the remaining
rows to $V$.  This is stable skewing followed by finite-row
specialization, as in Remark~\ref{rem:skew-specialization}.


\subsection{Simultaneous derivative embeddings}

Taylor expansion gives simultaneous maps behind the character
identities.  Derivative closure keeps every Taylor coefficient
inside the module.  Schur--Weyl duality then separates all
Littlewood--Richardson copies, so their contributions form an
actual direct sum.  This construction applies to any
derivative-closed polynomial module.

\begin{theorem}[The derivative coalgebra and Littlewood--Richardson embeddings]
\label{thm:derivative-coalgebra}
Let $V,W$ be finite-dimensional vector spaces over a field of
characteristic zero, let $G\subseteq\GL(W)$, and let
$\mathcal M\subseteq\Sym(V\otimes W)$ be a graded
$\GL(V)\times G$-submodule closed under all constant-coefficient
partial derivatives. Put
\[
 M_\lambda=\Hom_{\GL(V)}(\Schur_\lambda V,\mathcal M),
 \qquad M_\lambda=0\quad\text{if }\ell(\lambda)>\dim V.
\]
Then $\mathcal M$ is a subcoalgebra for the polynomial coproduct
$\Delta F(X,Y)=F(X+Y)$. For every pair of partitions $\tau,\rho$,
the coproduct induces a $G$-equivariant injection
\begin{equation}\label{eq:lr-coalgebra-injection}
 \bigoplus_\lambda M_\lambda^{\oplus c_{\tau,\rho}^{\lambda}}
       \lhook\joinrel\longrightarrow M_\tau\otimes M_\rho.
\end{equation}
All Littlewood--Richardson multiplicities occur independently.
The maps are compatible with inclusions of derivative-closed submodules.
\end{theorem}
\begin{proof}
Taylor expansion and derivative closure give
\[
 \Delta\mathcal M\subseteq
 [\Sym(V\otimes W)\otimes\mathcal M]
 \cap[\mathcal M\otimes\Sym(V\otimes W)]
 =\mathcal M\otimes\mathcal M.
\]
The counit is evaluation at zero. This proves the coalgebra assertion.

Put $a=|\tau|$ and $b=|\rho|$. Project the bidegree-$(a,b)$
coproduct to the row types $\tau$ and $\rho$ in its two factors.
The ambient Cauchy decomposition has column multiplicity
$\Schur_\lambda W$ in row type $\lambda$. Schur--Weyl duality
constructs this summand by the invariant pairing between two copies
of the Specht module $[\lambda]$.  Write
$L_{\tau,\rho}^{\lambda}$ for the multiplicity space of
$[\tau]\boxtimes[\rho]$ in its restriction to $S_a\times S_b$;
its dimension is $c_{\tau,\rho}^{\lambda}$.
The bidegree coproduct uses the restricted Specht pairing.
Over the reals, choose this pairing positive and invariant, so
its restriction to each isotypic component is nondegenerate.
On the two copies of $L_{\tau,\rho}^{\lambda}$ it is therefore a
perfect pairing.  In dual bases its tensor is the identity:
each row coefficient selects a distinct column coefficient.
We obtain exactly $c_{\tau,\rho}^{\lambda}$ independent ambient maps
\[
 \Schur_\lambda W\longrightarrow
                  \Schur_\tau W\otimes\Schur_\rho W.
\]
For different $\lambda$ their images lie in distinct isotypic
summands; for a fixed $\lambda$ they fill its independent
Littlewood--Richardson copies. Their direct sum is injective.
All relevant pairings and projectors can be chosen over $\Q$,
so the assertion holds over every characteristic-zero field.

The coalgebra inclusion places the projected coproduct inside
the row spaces tensored with $M_\tau\otimes M_\rho$. Taking its
independent row coefficients therefore restricts the ambient maps
to \eqref{eq:lr-coalgebra-injection}. The column group is used
only on the ambient polynomial ring; $\mathcal M$ need only be
$G$-stable. Fixed choices of ambient maps give compatibility
with inclusions.
\end{proof}

\begin{corollary}[Exterior derivative embeddings]
\label{thm:exterior-derivatives}
Under the hypotheses of Theorem~\ref{thm:derivative-coalgebra},
for every partition $\rho$ and $j\geq0$ there is a
$G$-equivariant injection
\begin{equation}\label{eq:exterior-injection}
 \bigoplus_{\lambda/\rho\text{ a vertical }j\text{-strip}}
       M_\lambda
 \lhook\joinrel\longrightarrow\Lambda^jW\otimes M_\rho.
\end{equation}
The maps arise from the exterior row component of the $j$th
derivatives and are compatible with inclusions.
\end{corollary}
\begin{proof}
Take $\tau=(1^j)$ in \eqref{eq:lr-coalgebra-injection} and use
the vertical Pieri rule and $M_{(1^j)}\subseteq\Lambda^jW$.
\end{proof}

\begin{corollary}[Derivative realization and exact hook defects]
\label{cor:hook-lower}
\label{cor:hook-realization}
\label{cor:hook-defects}
For $\mathcal M=\Harm$ or $\Pol$, let $M_\rho^{\mathcal M}$ denote
its stable row multiplicity space and put
$\mathcal A_\lambda^{\mathcal M}
=\Hom_{S_n}(\sgn,M_\lambda^{\mathcal M})$.
The exterior derivative maps give compatible isomorphisms
\begin{equation}\label{eq:hook-realization}
 \bigoplus_{\lambda/\rho\text{ a vertical }d\text{-strip}}
       \mathcal A_\lambda^{\mathcal M}
 \xrightarrow{\ \sim\ }
 \Hom_{S_n}(S^{(d+1,1^{n-d-1})},M_\rho^{\mathcal M}),
 \qquad 0\le d<n,
\end{equation}
and induce isomorphisms on the quotients for $\Pol\subseteq\Harm$.
Thus, writing $D_A=A_n^\Harm-A_n^\Pol$ and
$D_d=C_{(d+1,1^{n-d-1})}^\Harm-C_{(d+1,1^{n-d-1})}^\Pol$,
\begin{align}
 C_{(d+1,1^{n-d-1})}^{\mathcal M}
       &=e_d^\perp A_n^{\mathcal M},\label{eq:hook-lower}\\
 D_d&=e_d^\perp D_A.\label{eq:hook-defects}
\end{align}
Restricting source and target to $k$ rows gives
\[
 c_{\rho,(d+1,1^{n-d-1})}^{\mathcal M}
 \ge\sum_{\substack{\lambda/\rho\text{ vertical }d\text{-strip}\\
                    \ell(\lambda)\le k}}
            c_{\lambda,(1^n)}^{\mathcal M}.
\]
\end{corollary}
\begin{proof}
Apply Corollary~\ref{thm:exterior-derivatives} with $W=W_n$
and take $\Hom_{S_n}(\sgn,-)$.
Since $\sgn\otimes\Lambda^dW_n\cong S^{(d+1,1^{n-d-1})}$,
the target is the indicated hook multiplicity.
Theorems~\ref{thm:full-hook} and~\ref{thm:r2-operator-hook-transfer}
give equality of dimensions, so these compatible injections are
isomorphisms. They induce the quotient isomorphisms and both
character identities. Retaining only source partitions of length
at most $k$ gives the finite-row bound.
\end{proof}

Thus full alternating generation is equivalent to simultaneous
generation of every hook label. The derivative maps retain the
individual vertical-strip sources, and their quotients identify
every remaining hook defect.

\subsection{Two-column support and the alternating defect}
\label{sec:reconstruction}

After removing the linear moments, the labels form the standard
representation $W_n$ with its invariant inner product.  The
quadratic moments are its pairwise Gram products.  Two alternating
blocks whose total height exceeds $\dim W_n$ vanish modulo those
products.  Applied to the first two columns of a row partition,
this gives the support restriction used in the uniform generation proof.

\begin{lemma}[Two alternating blocks in the Gram ideal]
\label{lem:gram}
Let $E$ be a vector space of dimension $N$ over a
field of characteristic zero, equipped with a
nondegenerate symmetric bilinear form.
A multilinear polynomial on $E^{d+e}$ that is
alternating in the first $d$ arguments and in the
last $e$ arguments belongs to the ideal generated
by the inner products between the two blocks if $d+e>N$.
\end{lemma}

\begin{proof}
Identify $E$ with $E^*$ by the form. On the tensor product of exterior
algebras, using ordinary tensor signs, put
\[
 L=\sum_i\eps_i\otimes\eps_i,\qquad
 c=\sum_i\iota_i\otimes\iota_i.
\]
Invariantly the form and its inverse define these operators over $K$.
They satisfy $[c,L]=(N-d-e)\operatorname{id}$ in bidegree $(d,e)$.
Finite-dimensional $\mathfrak{sl}_2$ theory makes
\[
 L:\Lambda^{d-1}E\otimes\Lambda^{e-1}E
                 \longrightarrow\Lambda^dE\otimes\Lambda^eE
\]
surjective when $d+e>N$. Under the multilinear polynomial interpretation,
its image consists of alternating sums of a cross-block inner product
multiplied by a polynomial in the remaining arguments. This proves
membership in the stated ideal. An explicit finite right inverse is
recorded in \eqref{eq:gram-finite-certificate}.
\end{proof}

The next inequality is the admissibility condition for orthogonal
traceless tensors; see, for example, \cite{BGH}.  The Gram-ideal
argument proves it directly at the dimension $n-1$ relevant here.

\begin{theorem}[Two-column support]\label{thm:two-column}
For every coefficient of full diagonal harmonics,
\begin{equation}\label{eq:two-column}
 c_{\lambda,\mu}^{(n)}\ne0
       \quad\Longrightarrow\quad
       \lambda'_1+\lambda'_2\leq n-1.
\end{equation}
Here missing columns have length zero.
The same restriction holds in $\Pol$.
\end{theorem}

\begin{proof}
Quotient by the linear invariants first.  This
identifies the polynomial ring with
$\Sym(V\otimes W_n)$, where $V=K^k$ and
$\dim W_n=n-1$.
The quadratic invariants include all pairwise
inner products of the row vectors for the
nondegenerate form on $W_n$.
Thus the full coinvariant ring is a quotient of
the quotient by the Gram ideal.

The usual minor construction of the Cauchy
decomposition spans a row-$\lambda$ highest-weight
space by products of minors with column heights
$\lambda'_1,\lambda'_2,\ldots$.
If the sum of the first two heights exceeds
$n-1$, regard the vectors in these two minors
temporarily as independent.  Their product is
multilinear and alternating separately in the
two blocks.  Lemma~\ref{lem:gram} puts this product
in the Gram ideal.  Identifying the repeated
row vectors and multiplying by the remaining
minors keeps it in that ideal.
If only one column is present, a height greater
than $n-1$ is already excluded by the Cauchy
decomposition.
Every such highest-weight vector vanishes in
the quotient, proving \eqref{eq:two-column}.
The assertion for $\Pol$ follows by inclusion.
\end{proof}

\begin{corollary}[Support of the alternating defect]
\label{cor:r2-defect-support}
If the Schur coefficient of $s_\lambda$ in $D_A$ is nonzero, then
\begin{equation}\label{eq:r2-defect-support}
 \lambda_2\geq3,\qquad \lambda_3\geq2,\qquad
                 \lambda'_1+\lambda'_2\leq n-1.
\end{equation}
In particular, $\Pol$ contains every stable alternating component and
every stable hook-label component when $n\leq6$.
\end{corollary}

\begin{proof}
Haiman's two-row operator theorem gives $D_j(q_1,q_2)=0$ for every $j$.
If $\lambda_3\leq1$, delete one box from each nonempty row after the
second; the remaining partition has at most two rows, and the removed
boxes form a vertical strip.  Schur positivity of $D_A$ and the
vertical Pieri rule would then give a nonzero coefficient in some
$e_j^\perp D_A(q_1,q_2)$, a contradiction.

Every row-hook component of every label is generated, by
Proposition~\ref{prop:row-hook-generation}.  If $\lambda_2\leq2$, remove
the last box from every row of length two below the first.  The
remaining partition is a row hook.  Again positivity and
\eqref{eq:hook-defects} contradict generation of that
row-hook component.  The strip sizes used here are at most $n-1$,
since the row-length bound applies to the alternating label.

The final inequality is the two-column support theorem.  The first
two inequalities force the first two columns of $\lambda$ to have
at least three boxes each.  Thus $n\geq7$ for a nonzero defect.
Equation~\eqref{eq:hook-defects} gives the final assertion.
\end{proof}

\subsection{Finite-alphabet reconstruction}

The hook identities also recover stable alternating information
beyond the number of rows in which the hook characters are known.

\begin{theorem}[Reconstruction from finitely many alphabets]
\label{thm:reconstruction}
Let $\mathcal M=\Harm$ or $\Pol$, let $n\geq3$, and put
$k=\lfloor(n-1)/2\rfloor$.
The collection
\begin{equation}\label{eq:reconstruction-data}
 \left\{
 C_{(d+1,1^{n-d-1})}^{\mathcal M}(q_1,\ldots,q_k):
          0\leq d\leq n-1
 \right\}
\end{equation}
determines every Schur coefficient of the stable
alternating character $A_n^{\mathcal M}$.
\end{theorem}

\begin{proof}
Suppress the superscript $\mathcal M$, write
$A_n=\sum_\lambda a_\lambda s_\lambda$, and form
\[
 F(\bq,u)=\sum_{d=0}^{n-1}
 C_{(d+1,1^{n-d-1})}(\bq)u^d.
\]
The hook identity gives
\begin{equation}\label{eq:reconstruction-F}
 F(\bq,u)=\sum_\lambda a_\lambda
                         s_\lambda(\bq\mid u).
\end{equation}
If $\lambda_{k+1}\geq2$, then
$\lambda'_1+\lambda'_2\geq2k+2>n-1$.
Theorem~\ref{thm:two-column} excludes such a
partition.  Every supported partition therefore
satisfies $\lambda_{k+1}\leq1$.

At $u=0$, equation \eqref{eq:reconstruction-F}
determines all $a_\lambda$ with $\ell(\lambda)\leq k$,
because their Schur functions in $k$ variables
are linearly independent.  Subtract their entire
supersymmetric contributions from $F$.
Every remaining partition has the form
\[
 \lambda=(\lambda_1,\ldots,\lambda_k,1^r),
       \qquad r>0,\quad\lambda_k\geq1.
\]
The standard factorization of a supersymmetric
Schur function on the $(k,1)$ hook gives
\begin{equation}\label{eq:super-factorization}
 s_\lambda(\bq\mid u)
    =\prod_{i=1}^k(q_i+u)\,
       s_{(\lambda_1-1,\ldots,\lambda_k-1)}(\bq)\,u^r.
\end{equation}
It also follows directly from the vertical-strip
expansion \eqref{eq:superschur} and the Pieri rule.
Divide the residual polynomial by
$\prod_i(q_i+u)$.
The polynomials $s_\nu(\bq)u^r$ are linearly
independent, and the pair $(\nu,r)$ recovers
$\lambda=(\nu_1+1,\ldots,\nu_k+1,1^r)$.
This determines all remaining coefficients.
\end{proof}

Bergeron's reconstruction from hook information
\cite[Section~3.1]{Bergeron} is formulated under the corresponding
structural conjectures.  Theorem~\ref{thm:reconstruction} establishes
it for both the full and operator characters, with the number of
rows fixed by the two-column bound.  For seven labels, three-row
hook characters determine the entire stable alternating character.
